\documentclass[10pt,a4paper]{amsart}
\usepackage{amsmath,amssymb,amsthm,graphicx}
\usepackage{verbatim}
\newenvironment{reference}{\comment}{\endcomment}
%\newenvironment{reference}{}{}
\newenvironment{slogan}{\comment}{\endcomment}
\newenvironment{history}{\comment}{\endcomment}

% For commutative diagrams we use Xy-pic
\usepackage[all]{xy}

% We use 2cell for 2-commutative diagrams.
\xyoption{2cell}
% Theorem environments.
%
\theoremstyle{plain}
\newtheorem{theorem}[subsection]{Theorem}
\newtheorem{proposition}[subsection]{Proposition}
\newtheorem{lemma}[subsection]{Lemma}

\theoremstyle{definition}
\newtheorem{definition}[subsection]{Definition}
\newtheorem{example}[subsection]{Example}
\newtheorem{exercise}[subsection]{Exercise}
\newtheorem{situation}[subsection]{Situation}

\theoremstyle{remark}
\newtheorem{remark}[subsection]{Remark}
\newtheorem{remarks}[subsection]{Remarks}

\numberwithin{equation}{subsection}

% Macros
%
\def\lim{\mathop{\mathrm{lim}}\nolimits}
\def\colim{\mathop{\mathrm{colim}}\nolimits}
\def\Spec{\mathop{\mathrm{Spec}}}
\def\Hom{\mathop{\mathrm{Hom}}\nolimits}
\def\Ext{\mathop{\mathrm{Ext}}\nolimits}
\def\SheafHom{\mathop{\mathcal{H}\!\mathit{om}}\nolimits}
\def\SheafExt{\mathop{\mathcal{E}\!\mathit{xt}}\nolimits}
\def\Sch{\mathit{Sch}}
\def\Mor{\mathop{\mathrm{Mor}}\nolimits}
\def\Ob{\mathop{\mathrm{Ob}}\nolimits}
\def\Sh{\mathop{\mathit{Sh}}\nolimits}
\def\NL{\mathop{N\!L}\nolimits}
\def\CH{\mathop{\mathrm{CH}}\nolimits}
\def\proetale{{pro\text{-}\acute{e}tale}}
\def\etale{{\acute{e}tale}}
\def\QCoh{\mathit{QCoh}}
\def\Ker{\mathop{\mathrm{Ker}}}
\def\Im{\mathop{\mathrm{Im}}}
\def\Coker{\mathop{\mathrm{Coker}}}
\def\Coim{\mathop{\mathrm{Coim}}}

% Boxtimes
%
\DeclareMathSymbol{\boxtimes}{\mathbin}{AMSa}{"02}

%
% Macros for moduli stacks/spaces
%
\def\QCohstack{\mathcal{QC}\!\mathit{oh}}
\def\Cohstack{\mathcal{C}\!\mathit{oh}}
\def\Spacesstack{\mathcal{S}\!\mathit{paces}}
\def\Quotfunctor{\mathrm{Quot}}
\def\Hilbfunctor{\mathrm{Hilb}}
\def\Curvesstack{\mathcal{C}\!\mathit{urves}}
\def\Polarizedstack{\mathcal{P}\!\mathit{olarized}}
\def\Complexesstack{\mathcal{C}\!\mathit{omplexes}}
% \Pic is the operator that assigns to X its picard group, usage \Pic(X)
% \Picardstack_{X/B} denotes the Picard stack of X over B
% \Picardfunctor_{X/B} denotes the Picard functor of X over B
\def\Pic{\mathop{\mathrm{Pic}}\nolimits}
\def\Picardstack{\mathcal{P}\!\mathit{ic}}
\def\Picardfunctor{\mathrm{Pic}}
\def\Deformationcategory{\mathcal{D}\!\mathit{ef}}
\usepackage{cleveref}
\crefname{originaldefinition}{original definition}{original definitions}
\Crefname{originaldefinition}{Original definition}{Original definitions}
\makeatletter
\newlabel{definition-site-local}{{03DL}{source}{}{sourcepointer-1}{}}
\newlabel{definition-site-local@cref}{{[originaldefinition][0][]03DL}{[1][0][]source}{}{}{}}
\makeatother
\crefname{originalsection}{original section}{original sections}
\Crefname{originalsection}{Original section}{Original sections}
\makeatletter
\newlabel{section-intrinsic}{{03DG}{source}{}{sourcepointer-2}{}}
\newlabel{section-intrinsic@cref}{{[originalsection][0][]03DG}{[1][0][]source}{}{}{}}
\makeatother
\crefname{originaldefinition}{original definition}{original definitions}
\Crefname{originaldefinition}{Original definition}{Original definitions}
\makeatletter
\newlabel{definition-global}{{03DE}{source}{}{sourcepointer-3}{}}
\newlabel{definition-global@cref}{{[originaldefinition][0][]03DE}{[1][0][]source}{}{}{}}
\makeatother
\crefname{originalreference}{original reference}{original references}
\Crefname{originalreference}{Original reference}{Original references}
\makeatletter
\newlabel{sites-lemma-localize-special-cocontinuous}{{03CH}{source}{}{sourcepointer-4}{}}
\newlabel{sites-lemma-localize-special-cocontinuous@cref}{{[originalreference][0][]03CH}{[1][0][]source}{}{}{}}
\makeatother
\crefname{originallemma}{original lemma}{original lemmas}
\Crefname{originallemma}{Original lemma}{Original lemmas}
\makeatletter
\newlabel{lemma-global-pullback}{{03DF}{source}{}{sourcepointer-5}{}}
\newlabel{lemma-global-pullback@cref}{{[originallemma][0][]03DF}{[1][0][]source}{}{}{}}
\makeatother
\title{Intrinsicness of the local properties of modules on a site under passage to an equivalent topos}
\author{The Stacks Project Authors}
\date{stacks-project:sites-modules.tex:local-snapshot-20261012; GNU FDL 1.2 or later}
\begin{document}
\maketitle
\noindent\textit{Editorial note.} All delivered mathematical text is the author's own, copied verbatim from the pinned Stacks Project chapter file sites-modules.tex; no statement or proof step was rewritten, reordered or removed.\par\medskip
\noindent\textit{Editorial note.} The proof is long, so only the set-up, the reduction and the closing coherent-case argument were quoted for the roadmap; the counted span contains the entire author's proof (lines 3214-3332).\par\medskip
\noindent\textit{Editorial note.} The counted span cites the two definitions and a section of the chapter, recorded in assembly.external\_references with canonical tags as navigation pointers.\par\medskip
\noindent\textit{Allowed standard prerequisites (editorial).}\begin{itemize}
\item Special cocontinuous functors induce equivalences of topoi. Stacks Project, Sites, section on special cocontinuous functors (equivalence of topoi). Provides the functor g and its inverse used to transport modules and structure sheaves.
\item Intrinsicness of the global properties of modules on a site. Stacks Project, Modules on Sites, section on intrinsic properties (Definition of the local and global properties). Supplies the already-established intrinsicness of the global counterparts P\_g to which each local property is reduced.
\end{itemize}\medskip
\ifdefined\hypertarget\hypertarget{sourcepointer-1}{}\fi
\noindent\textit{Source reference.} Original definition 03DL: sources/stacks-chow/tags.txt, source lines 4424--4424. Cross-chapter standard prerequisite cited by the author inside the counted span; it is used as a pointer and its own proof is not part of this item.\par\medskip
\ifdefined\hypertarget\hypertarget{sourcepointer-2}{}\fi
\noindent\textit{Source reference.} Original section 03DG: sources/stacks-chow/tags.txt, source lines 4419--4419. Cross-chapter standard prerequisite cited by the author inside the counted span; it is used as a pointer and its own proof is not part of this item.\par\medskip
\ifdefined\hypertarget\hypertarget{sourcepointer-3}{}\fi
\noindent\textit{Source reference.} Original definition 03DE: sources/stacks-chow/tags.txt, source lines 4417--4417. Cross-chapter standard prerequisite cited by the author inside the counted span; it is used as a pointer and its own proof is not part of this item.\par\medskip
\ifdefined\hypertarget\hypertarget{sourcepointer-4}{}\fi
\noindent\textit{Source reference.} Original reference 03CH: sources/stacks-chow/tags.txt, source lines 4384--4384. Cross-chapter standard prerequisite cited by the author inside the counted span; it is used as a pointer and its own proof is not part of this item.\par\medskip
\ifdefined\hypertarget\hypertarget{sourcepointer-5}{}\fi
\noindent\textit{Source reference.} Original lemma 03DF: sources/stacks-chow/tags.txt, source lines 4418--4418. Cross-chapter standard prerequisite cited by the author inside the counted span; it is used as a pointer and its own proof is not part of this item.\par\medskip
\begin{lemma}
\label{lemma-special-locally-free}
Any of the properties (1) -- (8) of Definition \ref{definition-site-local}
is intrinsic (see discussion in Section \ref{section-intrinsic}).
\end{lemma}

\begin{proof}
Let $\mathcal{C}$, $\mathcal{D}$ be sites.
Let $u : \mathcal{C} \to \mathcal{D}$ be a special cocontinuous functor.
Let $\mathcal{O}$ be a sheaf of rings on $\mathcal{C}$.
Let $\mathcal{F}$ be a sheaf of $\mathcal{O}$-modules on $\mathcal{C}$.
Let $g : \Sh(\mathcal{C}) \to \Sh(\mathcal{D})$
be the equivalence of topoi associated to $u$.
Set $\mathcal{O}' = g_*\mathcal{O}$, and let
$g^\sharp : \mathcal{O}' \to g_*\mathcal{O}$ be the identity.
Finally, set $\mathcal{F}' = g_*\mathcal{F}$.
Let $\mathcal{P}_l$ be one of the properties (1) -- (7) listed in
Definition \ref{definition-site-local}.
(We will discuss the coherent case at the end of the proof.)
Let $\mathcal{P}_g$ denote the corresponding property listed in
Definition \ref{definition-global}. We have already seen that
$\mathcal{P}_g$ is intrinsic.
We have to show that
$\mathcal{P}_l(\mathcal{C}, \mathcal{O}, \mathcal{F})$
holds if and only if
$\mathcal{P}_l(\mathcal{D}, \mathcal{O}', \mathcal{F}')$
holds.

\medskip\noindent
Assume that $\mathcal{F}$ has $\mathcal{P}_l$.
Let $V$ be an object of $\mathcal{D}$.
One of the properties of a special cocontinuous functor is that there exists
a covering $\{u(U_i) \to V\}_{i \in I}$ in the site $\mathcal{D}$.
By assumption, for each $i$ there exists a covering
$\{U_{ij} \to U_i\}_{j \in J_i}$ in $\mathcal{C}$ such that
each restriction $\mathcal{F}|_{U_{ij}}$ is $\mathcal{P}_g$. By
Sites, Lemma \ref{sites-lemma-localize-special-cocontinuous}
we have commutative diagrams of ringed topoi
$$
\xymatrix{
(\Sh(\mathcal{C}/U_{ij}), \mathcal{O}_{U_{ij}}) \ar[r] \ar[d] &
(\Sh(\mathcal{C}), \mathcal{O}) \ar[d] \\
(\Sh(\mathcal{D}/u(U_{ij})), \mathcal{O}'_{u(U_{ij})}) \ar[r] &
(\Sh(\mathcal{D}), \mathcal{O}')
}
$$
where the vertical arrows are equivalences. Hence we conclude that
$\mathcal{F}'|_{u(U_{ij})}$ has property $\mathcal{P}_g$ also.
And moreover, $\{u(U_{ij}) \to V\}_{i \in I, j \in J_i}$ is a
covering of the site $\mathcal{D}$. Hence $\mathcal{F}'$ has
property $\mathcal{P}_l$.

\medskip\noindent
Assume that $\mathcal{F}'$ has $\mathcal{P}_l$.
Let $U$ be an object of $\mathcal{C}$.
By assumption, there exists a covering
$\{V_i \to u(U)\}_{i \in I}$ such that $\mathcal{F}'|_{V_i}$
has property $\mathcal{P}_g$. Because $u$ is cocontinuous we
can refine this covering by a family $\{u(U_j) \to u(U)\}_{j \in J}$
where $\{U_j \to U\}_{j \in J}$ is a covering in $\mathcal{C}$.
Say the refinement is given by $\alpha : J \to I$ and
$u(U_j) \to V_{\alpha(j)}$.
Restricting is transitive, i.e.,
$(\mathcal{F}'|_{V_{\alpha(j)}})|_{u(U_j)} = \mathcal{F}'|_{u(U_j)}$.
Hence by Lemma \ref{lemma-global-pullback} we see that
$\mathcal{F}'|_{u(U_j)}$ has property $\mathcal{P}_g$.
Hence the diagram
$$
\xymatrix{
(\Sh(\mathcal{C}/U_j), \mathcal{O}_{U_j}) \ar[r] \ar[d] &
(\Sh(\mathcal{C}), \mathcal{O}) \ar[d] \\
(\Sh(\mathcal{D}/u(U_j)), \mathcal{O}'_{u(U_j)})
\ar[r] &
(\Sh(\mathcal{D}), \mathcal{O}')
}
$$
where the vertical arrows are equivalences shows that $\mathcal{F}|_{U_j}$
has property $\mathcal{P}_g$ also. Thus $\mathcal{F}$ has
property $\mathcal{P}_l$ as desired.

\medskip\noindent
Finally, we prove the lemma in case
$\mathcal{P}_l = coherent$\footnote{The mechanics of this
are a bit awkward, and we suggest the reader skip this part of the proof.}.
Assume $\mathcal{F}$ is coherent. This implies that $\mathcal{F}$
is of finite type and hence $\mathcal{F}'$ is of finite type also by the
first part of the proof. Let $V$ be an object of $\mathcal{D}$ and let
$s_1, \ldots, s_n \in \mathcal{F}'(V)$. We have to show that the kernel
$\mathcal{K}'$ of
$\bigoplus_{j = 1, \ldots, n} \mathcal{O}_V \to \mathcal{F}'|_V$
is of finite type on $\mathcal{D}/V$. This means we have to show that
for any $V'/V$ there exists a covering $\{V'_i \to V'\}$ such that
$\mathcal{F}'|_{V'_i}$ is generated by finitely many sections.
Replacing $V$ by $V'$ (and restricting the sections $s_j$ to $V'$)
we reduce to the case where $V' = V$. Since $u$ is a special
cocontinuous functor, there exists a covering $\{u(U_i) \to V\}_{i \in I}$
in the site $\mathcal{D}$. Using the isomorphism of topoi
$\Sh(\mathcal{C}/U_i) = \Sh(\mathcal{D}/u(U_i))$
we see that $\mathcal{K}'|_{u(U_i)}$ corresponds to the kernel
$\mathcal{K}_i$ of a map
$\bigoplus_{j = 1, \ldots, n} \mathcal{O}_{U_i} \to \mathcal{F}|_{U_i}$.
Since $\mathcal{F}$ is coherent we see that $\mathcal{K}_i$
is of finite type. Hence we conclude (by the first part of the proof again)
that $\mathcal{K}|_{u(U_i)}$ is of finite type. Thus there exist coverings
$\{V_{il} \to u(U_i)\}$ such that $\mathcal{K}|_{V_{il}}$ is generated
by finitely many global sections. Since
$\{V_{il} \to V\}$ is a covering of $\mathcal{D}$ we conclude that
$\mathcal{K}$ is of finite type as desired.

\medskip\noindent
Assume $\mathcal{F}'$ is coherent. This implies that $\mathcal{F}'$
is of finite type and hence $\mathcal{F}$ is of finite type also by the
first part of the proof. Let $U$ be an object of $\mathcal{C}$, and let
$s_1, \ldots, s_n \in \mathcal{F}(U)$. We have to show that the kernel
$\mathcal{K}$ of
$\bigoplus_{j = 1, \ldots, n} \mathcal{O}_U \to \mathcal{F}|_U$
is of finite type on $\mathcal{C}/U$. Using the isomorphism of topoi
$\Sh(\mathcal{C}/U) = \Sh(\mathcal{D}/u(U))$
we see that $\mathcal{K}|_U$ corresponds to the kernel
$\mathcal{K}'$ of a map
$\bigoplus_{j = 1, \ldots, n} \mathcal{O}_{u(U)} \to \mathcal{F}'|_{u(U)}$.
As $\mathcal{F}'$ is coherent, we see that $\mathcal{K}'$ is of finite
type. Hence, by the first part of the proof again, we conclude
that $\mathcal{K}$ is of finite type.
\end{proof}

\end{document}
